% Archived .iscn concrete-syntax example, removed from main.tex (Section III,
% "The Extended i* Metamodel and the ISCN Oracle", sec:istar) on the grounds
% that Definition 1 (Goal marking, \label{def:goal-marking}) already gives
% .iscn its formal meaning (an assignment of expected values compared against
% mu_{Sigma_n}), so also carrying a full textual listing + walkthrough was
% redundant in the short (conference) version. Kept here for a future
% long/journal version where a worked example is still wanted alongside the
% definition.
%
% Original position in main.tex: right after the "MeetingScheduler"
% applied-example paragraph and the "Pending" paragraph, replacing the
% one-line summary that now reads "An \texttt{.iscn} file is the
% instance-level companion to \texttt{.istar}: it assigns the expected final
% marking for chosen actor instances, compared by Part D against the actual
% $\mu_{\Sigma_n}$ (Definition~\ref{def:goal-marking})."
%
% Requires (already loaded by main.tex): \usepackage{listings}

An \texttt{.iscn} file is the instance-level companion to \texttt{.istar}, as \texttt{.aol} (Section~\ref{sec:acl-language}) is to \texttt{.acl}: it adds no new goal, only the expected final marking for chosen actor instances. Listing~\ref{lst:mtg-iscn} gives \texttt{mtg\_expected.iscn} for the cast of Listing~\ref{lst:mtg-aol}: \texttt{p1}/\texttt{p2} (calendar) are expected \texttt{Fulfilled} on \texttt{CollectFromCalendar}, \texttt{p3} (phone) \texttt{Pending} there and \texttt{Fulfilled} on \texttt{HavePPCalled}, and both root goals, \texttt{init1.HaveMeetingOrganized} and \texttt{org1.HaveMeetingScheduled}, \texttt{Fulfilled}.

\begin{lstlisting}[basicstyle=\ttfamily\scriptsize,frame=single,breaklines=true,caption={Expected final marking in \texttt{mtg\_expected.iscn}.},label={lst:mtg-iscn}]
scenario MeetingSchedulerExpected for "mtg.istar" {
  instance init1 : Initiator;
  instance org1  : Organizer;
  instance sec1  : Secretary;
  instance p1, p2, p3 : Participant;

  assign p1.CollectFromCalendar = Fulfilled;
  assign p2.CollectFromCalendar = Fulfilled;
  assign p3.CollectFromCalendar = Pending;
  assign p3.HavePPCalled        = Fulfilled;

  assign init1.HaveMeetingOrganized = Fulfilled;
  assign org1.HaveMeetingScheduled  = Fulfilled;
}
\end{lstlisting}

Part D defines how each \texttt{assign} is compared against the actual $\mu_{\Sigma_n}$ (Definition~\ref{def:goal-marking}).
